\documentclass[11pt,reqno]{amsart}
\usepackage[T1]{fontenc}
\usepackage{lmodern}
\usepackage[margin=1.10in]{geometry}
\usepackage{microtype}
\usepackage{mathtools,amssymb}
\usepackage[hidelinks]{hyperref}
\hypersetup{pdftitle={Minimax entropy estimation beyond two derivatives},
  pdfauthor={Kaiwen Jin}}
\newtheorem{theorem}{Theorem}
\newtheorem{lemma}{Lemma}[section]
\newtheorem{proposition}[lemma]{Proposition}
\numberwithin{equation}{section}
\newcommand{\R}{\mathbb R}
\newcommand{\T}{\mathbb T}
\newcommand{\E}{\mathbb E}
\newcommand{\PP}{\mathbb P}
\newcommand{\one}{\mathbf 1}
\newcommand{\dd}{\,\mathrm d}
\newcommand{\Var}{\operatorname{Var}}
\newcommand{\osc}{\operatorname{osc}}
\newcommand{\cl}{\operatorname{clip}}
\newcommand{\Pois}{\operatorname{Pois}}
\newcommand{\cF}{\mathcal F}
\newcommand{\cV}{\mathcal V}
\title{Minimax entropy estimation beyond two derivatives}
\author{Kaiwen Jin}
\email{yui.ui.siki@gmail.com}
\subjclass[2020]{62G05, 62C20, 94A17}
\begin{document}
\begin{abstract}
We determine the minimax root-mean-square error for differential
entropy over zero-extended, cube-supported Lipschitz classes of
smoothness \(s>2\), in dimension \(d\ge3\) and with integrability
\(2\le p<\infty\). For sufficiently large fixed radius, the rate is
\((n\log n)^{-s/(s+d)}+n^{-1/2}\), which matches the known lower
bound. For \(s\ge d\) the upper bound follows from a joint likelihood
mixture over a finite Fourier net. For \(2<s<d\) we use signed
local polynomial projections. We bound how many cells have a projected value small in magnitude
relative to the mean on the enlarged cell. This permits summing
the bias of a recursive polynomial estimator across amplitude levels.
\end{abstract}
\maketitle
\section{Introduction}
For a density \(f\) on \([0,1]^d\), its entropy is
\(H(f)=-\int_{[0,1]^d} f\log f\), with \(0\log 0=0\) and natural
logarithms. Smoothness controls how well a density can be
approximated locally, but the logarithm near zero complicates the
passage from density approximation to entropy estimation.
Han, Jiao, Weissman and Wu \cite[Theorem~1, p.~3229]{HJWW} obtained the sharp rate for
smoothness at most two and proved a lower bound that also applies
at higher smoothness. They discuss the higher-smoothness upper
bound in Section 4.2 of their author manuscript \cite{HJWWv4}.

Here we use the centered differences of order \(r=\lceil s\rceil\).
For \(u\in\R^d\) and \(t>0\), set
\[
 \Delta_u^r f(x)=\sum_{j=0}^r(-1)^{r-j}\binom rj
    f\bigl(x+(j-r/2)u\bigr),\qquad
 \omega_r(f,t)_p=\sup_{|u|\le t}\|\Delta_u^r f\|_{L^p(\R^d)} .
\]
The shift norm is Euclidean. Write \(\cF_{s,p,d}(L)\) for the
nonnegative probability densities supported on the unit cube with
\[
 \|f\|_{L^p(\R^d)}+\sup_{t>0}t^{-s}\omega_r(f,t)_p\le L .
\]
The differences act on the zero extension. This condition includes
behavior at the boundary and allows densities to vanish on sets of
positive measure.

For \(n\) independent observations with density \(f\), let
\[
 R_n=\inf_{\widehat H}\sup_{f\in\cF_{s,p,d}(L)}
          \bigl(\E_f|\widehat H-H(f)|^2\bigr)^{1/2},
\]
where the infimum is over measurable estimators.
\begin{theorem}\label{thm:main}
Fix an integer \(d\ge3\), a real \(s>2\), and \(2\le p<\infty\).
There is an \(L_0(d,s,p)>1\) such that for each fixed \(L\ge L_0\)
there are constants \(0<c\le C<\infty\) and \(n_0\ge2\), depending
only on \(d,s,p,L\), for which
\[
 c\bigl((n\log n)^{-s/(s+d)}+n^{-1/2}\bigr)
 \le R_n\le
 C\bigl((n\log n)^{-s/(s+d)}+n^{-1/2}\bigr)
 \qquad(n\ge n_0).
\]
\end{theorem}
The estimators establish the rate over the displayed class. Our
construction is for measurable estimation and gives no bound on
computational cost.

For the lower bound, the final lower-bound clause of Theorem 1 in
\cite{HJWWv4} allows all \(s>0\) and \(1\le p<\infty\).
Equations (1)--(4) there use the same zero-extended, cube-supported
density class, and (5) uses the same loss. The shifts \(te\), \(|e|\le1\),
are exactly the shifts \(|u|\le t\) in the definition above. The source lower bound is
\[
 R_n\ \ge\ c_{s,p,d}
 \bigl(L^{d/(s+d)}(n\log n)^{-s/(s+d)}+(\log L)n^{-1/2}\bigr)
\]
for \(L\) above a parameter-dependent threshold and
\(L\le(n\log n)^{s/d}\). The upper restriction on \(L\) holds
eventually at each fixed radius, so this yields the claimed lower bound.

There are two parts to the upper bound. A likelihood mixture
attains the parametric rate when \(s\ge d\). For \(2<s<d\), we
approximate \(f\) on enlarged mesh cells by polynomials and use
the odd extension \(z\log|z|\). The estimates on a cell depend
on its count and on the geometry of its projected shape. The
localization estimate in Section~\ref{sec:geometry} controls the
number of cells on which that shape approaches zero relative to
its mean. Section~\ref{sec:angular} constructs polynomial
estimators with a small power loss confined to these cells.
Sections~\ref{sec:high} and~\ref{sec:risk} give the large-count
variance bound and assemble the fixed-sample estimator.

\section{Fourier estimates and the range \texorpdfstring{\(s\ge d\)}{s >= d}}
\label{sec:large-s}

Identify the torus \(\T^d\) with the half-open unit cube and use
its volume-one measure. Periodize the zero extension of \(f\).
It agrees with \(f\) almost everywhere on the cube and has
\(\|f\|_2\le L\) and mean one. We use \(f\) also for the periodized
function from here on. Constants \(c,C\) may change between
occurrences and depend on the fixed class parameters. Further
fixed parameters in the construction are chosen in terms of those
parameters.

For \(|u|\le\sqrt d\), the supports of the terms of the
whole-space difference lie in a fixed box. Hölder's inequality in
that box bounds its \(L^2\) norm by \(C|u|^s\). The periodization
of the difference is the torus difference. Only a bounded number
of terms can overlap at each point in this periodization, so
Cauchy--Schwarz gives \(\|\Delta_u^r f\|_{2,\T^d}\le C|u|^s\).
On a Fourier mode \(k\in\mathbb Z^d\), the centered difference
has multiplier \((2i\sin(\pi k\cdot u))^r\). For \(M\ge1\), average
the squared norm over \(u\in[-M^{-1},M^{-1}]^d\). At any
frequency \(|k|_\infty>M\), one coordinate of \(k\cdot u\)
ranges over an interval longer than two as that coordinate of \(u\)
varies. The average of \(|2\sin(\pi k\cdot u)|^{2r}\) is
bounded below by a positive constant, by retaining the complete
periods in that interval. Parseval and Tonelli give
\begin{equation}\label{eq:tail}
   \sum_{|k|_\infty>M} |\widehat f(k)|^2\le C M^{-2s}.
\end{equation}
All estimates use the original difference order, including \(r=s\)
at integer \(s\). In the real Fourier shell
\(2^{j-1}<|k|_\infty\le2^j\) the dimension is \(O(2^{jd})\)
and the \(L^2\) norm is \(O(2^{-js})\). Use the shell
\(0<|k|_\infty\le1\) when \(j=0\).

To construct the mixture, fix \(M=2^J\), \(\zeta=M^{-s}\).
In shell \(j\le J\) choose a deterministic
\(\zeta 2^{-(J-j)}\)-net of its coefficient ball of radius
\(C2^{-js}\). The logarithm of the covering cardinality of a
radius-\(R\) ball in dimension \(m\) at resolution \(\rho\)
is at most \(m\log(1+2R/\rho)\), by comparing disjoint
radius-\(\rho/2\) balls centered at a maximal separated set.
Sum one choice from each shell and adjoin the constant one to
obtain the family \(\mathcal G_M\). Its members have uniformly
bounded \(L^2\) norms. By orthogonality and \eqref{eq:tail},
each \(f\) has a representative \(g\in\mathcal G_M\)
with \(\|g-f\|_2\le C\zeta\). The covering count is
\[
 \log|\mathcal G_M|\le C\sum_{j=0}^J2^{jd}(1+J-j)\le C M^d .
\]
For each \(g\), put \(q_g=(g_++\zeta)/Z_g\),
where \(Z_g=\int(g_++\zeta)\) is in \([1,C]\). For a representative
\(g\) of \(f\), truncation reduces the distance to \(f\),
\(|Z_g-1|\le \|g_++\zeta-f\|_1\le C\zeta\), and hence
\[
 \|q_g-f\|_2\le C\zeta,\qquad
 \chi^2(f,q_g):=\int (f-q_g)^2/q_g\le C\zeta .
\]
The floor \(q_g\ge c\zeta\) applies to the likelihoods used
in the estimator.

We use the finite-net likelihood comparison of Yang and Barron
\cite[Section~2, equation (2)]{YangBarron} as a pointwise
comparison and bound both signs in \(L^2\). Let
\(Q_n=|\mathcal G_M|^{-1}\sum_{g\in\mathcal G_M} q_g^{\otimes n}\)
and use \(\widehat H=-n^{-1}\log Q_n(X_1,\ldots,X_n)\).
It is finite and measurable on the product cube. If
\(B=n^{-1}\log(f^{\otimes n}/Q_n)(X_1,\ldots,X_n)\), then
\(\E_f e^{-nB}\le1\), by integrating the density \(Q_n\)
over the set where \(f^{\otimes n}>0\). Thus
\(\PP_f(B<-t)\le e^{-nt}\) and \(\|B_-\|_2\le \sqrt{2}/n\).
For the representative \(q=q_g\), the corresponding term
in the mixture bounds \(B\) above by
\[
 n^{-1}\log|\mathcal G_M|
 +n^{-1}\sum_{i=1}^n\log(f(X_i)/q(X_i)).
\]
The last average has mean in \([0,\chi^2(f,q)]\) by Jensen
and \(\log z\le z-1\). Its variance is at most \(C\zeta/n\)
since \(z(\log z)^2\le C(z-1)^2\) for \(z\ge0\), by
continuity at \(0,1\) and the bound at infinity. Consequently
\(\|B_+\|_2\le CM^d/n+C\zeta+C\sqrt{\zeta/n}\).
The empirical mean of \(-\log f(X_i)\) has mean \(H(f)\)
and standard deviation at most \(Cn^{-1/2}\), since
\(z(\log z)^2\le C(1+z^2)\). Adding \(B\) to that mean yields
\[
 \|\widehat H-H(f)\|_{L^2(\PP_f)}
 \le C\bigl(n^{-1/2}+M^{-s}+M^d/n\bigr).
\]
Take \(M\asymp n^{1/(s+d)}\) dyadic to get \(Cn^{-1/2}\)
for every fixed \(s\ge d\), including \(s=d\).

\section{Local projections and amplitude levels}\label{sec:geometry}

For the remainder of the proof suppose \(2<s<d\).
Partition \(\T^d\) into half-open dyadic cubes \(I\) of side
\(h\le1/8\). For each \(I\), enlarge its Euclidean lift by a
factor of three about its center and denote the resulting cube
by \(J_I\), with volume \(v=(3h)^d\). Its image \(J_I^{\rm phys}\)
on the torus is Borel, and the quotient restricted to the lift
is injective and measure preserving. Integrals and restrictions
to \(J_I\) refer to this lift with the periodized density. The
physical supports have overlap at most \(5^d\). Indeed, the center
of any support containing a point lies within circular distance
\(3h/2\) in each coordinate on a grid of spacing \(h\). Two
intersecting supports have centers within circular distance \(3h\)
in every coordinate. Thus the graph joining intersecting physical
supports has maximum degree at most \(7^d-1\).

Use real polynomials of total degree at most \(r=\lceil s\rceil\).
Let \(Q_I\) be the ordinary \(L^2\) projection of \(f\) on \(J_I\)
and \(e_I=\|f-Q_I\|_{2,J_I}\). The kernel on the volume-one cube
\(\Omega=[0,1)^d\) is \(K_\Omega(x,y)=\sum_i\phi_i(x)\phi_i(y)\)
for a real orthonormal basis. It is uniformly bounded and
reproduces the polynomial space. Rescaling the kernel shows that
\(Q_I\) has mean \(b_I=P_I/v\) and supremum norm at most \(C_V b_I\),
where \(P_I=\int_{J_I} f\). Write \(\mu_I=\int_I f\).
Under the affine map \(T_I\) from \(\Omega\) into the torus
through \(J_I\), the central cube \(W=[1/3,2/3)^d\)
maps to \(I\). The volume Jacobian of \(T_I\) is \(v\),
and \(vb_I=P_I\). If \(P_I>0\), set
\[
 w_I=f\circ T_I/b_I,\qquad q_I=Q_I\circ T_I/b_I,\qquad
 E_I=\|w_I-q_I\|_{2,\Omega}.
\]
Then \(w_I\) is a probability density, \(q_I\) is its ordinary
projection, \(\mu_I=P_I\int_W w_I\), and \(P_I E_I=\sqrt v\,e_I\).
If \(P_I=0\) the density and projection vanish almost everywhere on
\(J_I\), and products \(P_I E_I\) are defined to be zero.

We record the projection errors needed below. For every dyadic
\(h\le1/8\) and each dyadic \(\ell\in[h,1]\), let \(P_\ell\) be the
ordinary piecewise projection on disjoint cubes of side \(\ell\).
With gradients taken separately in the cubes,
\begin{equation}\label{eq:projection}
 \sum_I\sqrt v\,e_I\le C h^s,\qquad
 \|\nabla(f-P_\ell)\|_{2,\mathrm{piecewise}}\le C\ell^{s-1}.
\end{equation}
To see this, at side \(\ell\) split \(f\) at frequency of order
\(\ell^{-1}\). The tail and its gradient are \(O(\ell^s)\) and
\(O(\ell^{s-1})\) in \(L^2\). Indeed \eqref{eq:tail}
gives squared shell bounds \(C2^{-2js}\), and the
squared gradient tail is bounded by
\(C\sum_{2^j\gtrsim \ell^{-1}} 2^{-2j(s-1)}\).
In particular \(f\in W^{1,2}(\T^d)\). For the smooth
part \(g\), on any cube of side \(\ell\) choose the degree-\(r\)
polynomial \(P\) for which the derivatives of \(g-P\) through
order \(r\) have mean zero. Its coefficients are obtained in
descending order of total degree using the nonzero constant
derivative of the matching monomial. Poincaré's inequality on
each derivative with mean zero, iterated, gives
\[
 \|D^j(g-P)\|_2\le C\ell^{r+1-j}\|D^{r+1}g\|_2,\qquad j=0,1 ,
\]
on that cube, with the norms summing over derivatives of each order.
If \(\Pi\) is the local orthogonal projection, then
\(g-\Pi g=(g-P)+\Pi(P-g)\). Apply \(L^2\) contraction and
\(\|\nabla \Pi(P-g)\|_2\le C\ell^{-1}\|\Pi(P-g)\|_2\)
to retain the estimates for \(j=0,1\). The inverse inequality
follows by rescaling and finite-dimensional norm equivalence.
It also bounds the gradient of the projected tail. Over the torus,
\(\|D^{r+1}g\|_2\le C\ell^{s-r-1}\) by summing shells, since
\(r+1>s\). The argument works on the bounded-overlap family
\(J_I\) at side \(3h\) as well, giving
\(\sum_I e_I^2\le C h^{2s}\). Cauchy--Schwarz and
\(\sum_I v=3^d\) give the first bound in \eqref{eq:projection}.
Every estimate holds uniformly over the class on this mesh range.

\begin{lemma}\label{lem:sublevel}
For fixed degree bound \(r\), there are \(\beta>0\) and \(C<\infty\)
such that any real polynomial \(q\) of degree at most \(r\) on
\(\Omega\) with \(\|q\|_1\ge1\) satisfies
\[
 |\{x\in\Omega:|q(x)|\le z\}|\le C z^\beta
 \quad(z>0),\qquad \|\log^-|q|\|_{2,\Omega}\le C .
\]
\end{lemma}
\begin{proof}
In one dimension, if the sublevel set has length \(\lambda>0\),
choose \(r+1\) points from it with successive pairwise distances at
least \(\lambda/(4(r+1))\) by excluding intervals around the
points already chosen. Lagrange interpolation yields
\(\|q\|_\infty\le C z\lambda^{-r}\). Constant polynomials satisfy
the claim directly. For induction on dimension, write
\(q(x_1,x')=\sum_{i=0}^r x_1^i q_i(x')\).
Some coefficient has \(L^1\) norm at least \(1/(r+1)\). Outside
a set of measure \(O(z^{\beta'/2})\), its magnitude is at least
\(\sqrt z\) for \(0<z<1\), by the induction bound with exponent
\(\beta'\), scaled by its norm. A slice with such a coefficient
has supremum at least \(c\sqrt z\) by finite-dimensional norm
equivalence. The interval bound and Fubini prove the induction.
The zero set is null, and the tail of \(\log^-|q|\) is bounded
by \(C e^{-\beta t}\) for \(t\ge0\). Integrating the squared
tail formula gives the second bound.
\end{proof}

Put \(\psi(z)=z\log|z|\) for \(z\ne0\) and \(\psi(0)=0\).
Our deterministic surrogate on cell \(I\) is
\begin{equation}\label{eq:surrogate}
 F_I=P_I\int_W\psi(q_I)+\mu_I\log b_I
\end{equation}
for positive mass, and zero otherwise. For \(u\ge0\) and \(z\ne0\)
the divided difference of \(\psi\) in absolute value is at most
\[
 C+\log^+(u+|z|)+\log^-|z|.
\]
For \(z>0\) it is \(\log z + x\log x/(x-1)\) when \(x=u/z<1\),
and \(\log u+\log x/(x-1)\) when \(x>1\). The logarithmic ratios
here belong to \([0,1]\), including the endpoint limits. For \(z<0\)
the quotient is \(\log(u+|z|)+x\log x+(1-x)\log(1-x)\),
where \(x=u/(u+|z|)\). The last two terms have bounded sum.
As \(w_I\) has mass one and \(q_I\) has bounded \(L^1\) norm and
mean one, \((\log^+ t)^2\le C t\) and Lemma~\ref{lem:sublevel}
bound the divided difference in \(L^2(W)\) uniformly. Thus
Cauchy--Schwarz and the change of variables to \(W\) yield
\begin{equation}\label{eq:transfer}
 \sum_I\left|\int_I f\log f-F_I\right|
 \le C\sum_I P_I E_I\le C h^s .
\end{equation}
The scale term \(\mu_I\log b_I\) is common to the integral and
the surrogate, since it is formed with the mass of \(w_I\) on \(W\).

For \(P_I>0\) set \(\nu_I=\one_{\{\inf_W|q_I|<1/2\}}\),
and take \(\nu_I=0\) at zero mass.
This indicator records when the infimum of the absolute projected value
on \(I\) is below half the mean on \(J_I\).
Set \(\tau=h^s\) and \(\omega=2/(2s-1)\).
\begin{lemma}\label{lem:bands}
For \(B\ge\tau\), let \(\mathcal V_B\) contain the cells with
\(\nu_I=1\) and \(B/2<b_I\le B\). At \(\xi=B/\tau\),
\[
 h^d\#\mathcal V_B\le
 C\left(\xi^{-\omega}+h+
         B^{-1}\sum_{B/2<b_I\le B}\sqrt v\,e_I\right).
\]
\end{lemma}
\begin{proof}
Take \(U=f\wedge (C_0B)\) with \(C_0\ge \max(1,2C_V)\).
Let \(\Phi(y)=y^2\) on \(0\le y\le1\) and \(\Phi(y)=2y-1\)
on \(y>1\). On each \(J_I\),
\[
 \|U-\overline U_{J_I}\|_2^2
 \le C B^2\sum_{j=1}^d\int_{J_I}\Phi(h|\partial_j U|/B).
\]
On a coordinate slice, a pairwise difference is bounded by both
\(C_0 B\) and the integral of the absolute derivative on the whole
slice. Its square is bounded by \(CB^2\Phi\) of the latter integral
divided by \(B\). Use Jensen along the slice, whose length is \(3h\),
and average the pairwise differences. In the product cube, telescope
in coordinates and sum their squares, or use the coordinate
averaging contractions, to obtain the displayed inequality.

On each macro cube of side \(\ell\), fix a slice in direction \(j\)
and write \(g=\partial_j f\), \(k=\partial_j P_\ell\), \(M=B/h\).
The scalar bound
\[
 \Phi(|g|/M)\le
 C\bigl(\cl(k/M,[-1,1])\,g/M+(g-k)^2/M^2\bigr)
\]
holds for all real \(g,k\). For \(|k|\le M\), completing the
square controls \(g^2/M^2\). For \(k/M\ge1\), the right side
without the constant is \(y+(y-k/M)^2\) with \(y=g/M\).
For \(y\ge1/2\) use the linear term, for \(0\le y<1/2\)
use the squared term, and for \(y<0\) use
\(-|y|+(|y|+1)^2\). Reversing both signs covers \(k/M\le-1\).
Apply the inequality on \(\{f<C_0 B\}\), where
\(\partial_j U=g\). Elsewhere \(\partial_j U=0\) almost
everywhere by the Sobolev chain rule for a truncation.
The integral of \(\cl(k/M,[-1,1])\,\partial_j U/M\) on the
slice is at most \(C B/M\). Indeed the clipped polynomial
has bounded variation uniformly, from its finitely many
monotonicity intervals, and integration by parts also bounds the
two endpoint contributions. Sobolev slicing justifies the
one-dimensional integrals. The cross-sectional volume of each
macro cube is \(\ell^{d-1}\) and there are \(\ell^{-d}\) cubes,
so the linear term contributes \(O(h/\ell)\). The squared
term contributes at most \(M^{-2}\|\nabla(f-P_\ell)\|_2^2\)
up to a constant. Hence \eqref{eq:projection} gives
\begin{equation}\label{eq:huber}
 \sum_j\int_{\T^d}\Phi(h|\partial_j U|/B)
 \le C\left(h/\ell+h^2\ell^{2s-2}/B^2\right).
\end{equation}

In the band under study, \(Q_I\le C_VB\), so
\(\|f-U\|_{2,J_I}\le e_I\) and \(\|U-Q_I\|_{2,J_I}\le2e_I\).
The centered functions \(U-\overline U_{J_I}\) and \(Q_I-b_I\)
are bounded by \(CB\). Comparing squares and using
Cauchy--Schwarz gives
\(\|Q_I-b_I\|_{2,J_I}^2\le
\|U-\overline U_{J_I}\|_2^2+CB\sqrt v\,e_I\).
For \(\nu_I=1\) the supremum of \(|Q_I-b_I|\) is at least
\(b_I/2\), and by norm equivalence its squared \(L^2(J_I)\)
norm is at least \(cB^2 v\). Sum, use the bounded overlap
and \eqref{eq:huber}, then divide by \(B^2\). This proves the
claimed estimate with \(h/\ell+h^2\ell^{2s-2}/B^2\)
in place of \(\xi^{-\omega}+h\). When \(h\xi^\omega\le1\),
take a dyadic \(\ell\) within a factor of two of \(h\xi^\omega\)
in \([h,1]\). The two terms are \(O(\xi^{-\omega})\).
When \(h\xi^\omega>1\), use \(\ell=1\) and
\(h^{2-2s}\xi^{-2}\le h\).
\end{proof}

\section{A polynomial estimator for the cell functional}
\label{sec:angular}

We construct an estimator of the signed entropy of an ordinary
projection \(q\) on \(\Omega\), evaluated on \(W\).
In this section, \(w\ge0\) is in \(L^2(\Omega)\), with mass
\(p_0\), and \(e=\|w-q\|_2\). Fix an integer \(J\ge1\) and set
\(\alpha=3/4\), \(\epsilon=\alpha^J/2\).
Write \(K=\lfloor c_K\log n\rfloor\), with \(c_K>0\)
sufficiently small. We use two observation models.
\begin{enumerate}
\item For a Poisson correction cloud with intensity \(m w(x)\dd x\)
  on \(\Omega\), with \(m\asymp n\), assume \(p_0\le A\), where
  \(A=C_A\log n/n\). The cloud is the restriction of a parent
  sample of mean \(m\). Use \(a_*=A\) and \(r_*=A/K^2\).
\item For \(N\) observations from the probability density \(w\),
  with \(C_*\log n\le N\le n\), use an interval pilot and
  an independent correction fraction, each of size comparable
  to \(N\). Here \(p_0=1\), \(a_*=1\),
  \(t=\log n/N\) and \(r_*=t/K^2\).
\end{enumerate}
For fixed \(J\) and any \(\eta>0\), we will estimate
\(\int_W [\psi(q)-q\log a_*]\) with absolute bias at most
\begin{equation}\label{eq:angular-bias}
 C\left[e+r_*\left(1+(p_0/r_*)^\epsilon
                      \one_{\{\inf_W|q|<p_0/2\}}\right)\right].
\end{equation}
The second moment bounds will be \(C n^\eta A^2\log^2 n\) in
the Poisson model and \(C n^\eta\log^2 n\) in the fixed-count
model. In the latter model the bounds condition on the interval
pilot on an event with probability at least \(1-n^{-M}\)
for any prescribed \(M\). We describe the event and the statistic
on arbitrary pilot histories.

\subsection{Regional scores and angular intervals}

For a Borel set \(S\subset\Omega\) with \(v_S=|S|>0\),
we need a polynomial reproducing kernel on \(S\) with bounded
observation scores. On the polynomial space of dimension \(D\)
use the norm \(v_S^{-1}\int_S|P|\), which is definite by
Lemma~\ref{lem:sublevel}. Maximize the absolute determinant of
\(D\) vectors in its compact unit ball. The maximizing vectors
\(\varphi_i\) have norm one, since the maximum is positive.
Replacing one by another vector in the ball shows that its
coordinate functional has norm at most one.
Extend these functionals by Hahn--Banach to \(L^1\) of normalized
volume on \(S\). Duality gives representatives
\(\sigma_i\) bounded by one, with
\(v_S^{-1}\int_S\sigma_i\varphi_j=\delta_{ij}\).
Use finite Borel versions, clipped to \([-1,1]\) on \(S\)
and extended by zero off \(S\). Changing a Lebesgue null set
preserves these pairings and the score means in both models.

Let
\[
 B_S=\sum_i|\varphi_i|,\qquad
 \ell_S=B_S/v_S,\qquad
 K_S(x,y)=v_S^{-1}\one_S(y)\sum_i\varphi_i(x)\sigma_i(y).
\]
Since the constant one has all coordinates of magnitude at most
one, \(B_S\ge1\). Also \(\int_S B_S=Dv_S\).
The kernel reproduces every polynomial in the space and obeys
\(|K_S(x,y)|\le\ell_S(x)\one_S(y)\). Define
\(q_S=\int K_S(\cdot,y)w(y)\dd y\), \(c_S=q_S/\ell_S\)
and \(p_S=\int_S w\). Subtracting the reproduction of \(q\)
inside the integral and using Cauchy--Schwarz gives
\begin{equation}\label{eq:regional}
 |c_S|\le p_S,\qquad \int_S\ell_S=D,\qquad
 \|q_S-q\|_{1,S}\le D\sqrt{v_S}\,e .
\end{equation}
The actual scores \(g_{S,0}=\one_S\) and
\(g_{S,i}=\one_S(1+\sigma_i)/2\) belong to \([0,1]\)
everywhere. Their means against \(w\) are in \([0,p_S]\).
Using \(2\int w g_{S,i}-p_S\) for the signed coordinate
makes \(c_S(x)\) a linear form in the means with total absolute
coefficient bounded by three. For the ordinary root
\(S_0=\Omega\) use \(q_0=q\) and a constant envelope
\(\ell_0\ge\max(1,\|K_\Omega\|_\infty)\). The ordinary basis
scores are affine rescalings of the bounded basis polynomials
to \([0,1]\), together with the mass indicator. Their linear
form for \(c_0=q/\ell_0\) has bounded coefficient sum, and
\(|c_0|\le p_0\). The first kernel uses \(\Omega\) and its
first cut acts on \(X_0=W\). Each later kernel uses the
remaining part of that evaluation set.

In the Poisson model each score mean has interval \([0,A]\).
In the other model an empirical mean \(\widehat u\) gives
\[
 [l,l+B]=[(\widehat u-\Delta)_+,\widehat u+\Delta],
 \qquad \Delta=L_1(\sqrt{t\widehat u}+t).
\]
These intervals always have width at least \(L_1t\).
For a true mean \(u\), Bernstein's inequality uses
\(0\le g\le1\) and \(\Var(g)\le u\) to give
\(|\widehat u-u|\le C(\sqrt{tu}+t)\) at any fixed
inverse-polynomial failure probability. This implies
\(\sqrt{t\widehat u}+t\asymp\sqrt{tu}+t\), by separating
\(u\le C't\) from \(u>C't\) with \(C'\) sufficiently large.
Thus, with \(L_1\) large enough,
\begin{equation}\label{eq:interval}
 u\in[l,l+B],\qquad B\asymp\sqrt{tu}+t
\end{equation}
on that event. The fixed dictionaries of scores used below
have polynomial size in \(n\), so this bound holds simultaneously
with probability at least \(1-n^{-M}\) by a union bound.
Every interval lies in a fixed bounded nonnegative range,
including on failing histories.

Write \(u=l+B(1+\cos\theta)/2\) for one interval of width \(B\).
The derivative satisfies \(|u'(\theta)|^2=(u-l)(l+B-u)\).
In either model, valid intervals give the increment bound
\begin{equation}\label{eq:increment}
 |u(\theta+z)-u(\theta)|
 \le C\left(\sqrt{r_*p_S}\,K|z|+r_*(K|z|)^2\right)
\end{equation}
for a score supported on \(S\). In the Poisson model
\(B=A\), \(l=0\), and the first derivative has magnitude at most
\(\sqrt{Au}\) while the second derivative is \(O(A)\).
For the fixed-count model and \(u\ge t\), use
\(B|z|\le C\sqrt{tu}|z|\).
For \(u<t\) use \(B\le Ct\), first derivative bounded by
\(\sqrt{B(u-l)}\) and second derivative \(O(B)\).
The same increment estimate with a different constant applies
to \(c_S(x)\), uniformly in \(x\) and with a norm of the angle
vector in place of \(|z|\). Set \(D_S=r_*+\sqrt{r_*p_S}\).

\subsection{Polynomial selection and approximation}

On the angle circle, use a nonnegative even trigonometric
polynomial density of degree at most \(K\). For any fixed
integer \(b\ge1\) and large \(K\), we can take the normalized
\(2b\)-th power of the modulus of
\(\sum_{k=0}^{\lfloor K/b\rfloor}e^{ikz}\).
Its bound on \([-\pi,\pi]\) is
\(C_b K(1+K|z|)^{-2b}\). The geometric sum is bounded by
\(C\min(K,|z|^{-1})\) and has magnitude at least \(c_b K\)
on an interval of length \(c_b/K\), which yield the bound.
To define path probabilities at the true means, shift the
angles independently by this density, transform back to the
score intervals, and round on a nonnegative mesh of spacing at
most \(n^{-3}\). The statistic uses the polynomial path
probabilities in its coefficients, as shown below.
This size is \(O(r_*)\) throughout our count ranges.

For any rounded vector its probability, as a function of true
angles, is convolution of the product shift density with the
indicator of a bin. The bin is invariant under sign change in
each angle, so the probability is an even trigonometric
polynomial in each angle, hence a polynomial in the normalized
means \((u-l)/B\), of degree \(O(K)\). The sum of the
coefficient norms over all bins is at most \(C^K\).
Indeed the Fourier coefficients of the shift density, using
normalized Haar measure, have modulus at most one.
Each Fourier coefficient of a bin indicator is bounded by
its Haar measure, and those bin measures sum to one.
The number of Fourier indices is polynomial in \(K\).
The monomial coefficient norm of
\(\cos(k\theta)=T_k(2(u-l)/B-1)\) is at most \(C^k\)
by the Chebyshev recurrence. Multiplying over the fixed
number of angles gives the summed bound. Here and below the
coefficient norm is the sum of absolute monomial coefficients
in the normalized means.

Use a separate signed filter for approximation. In each score
angle its Fourier multiplier is \(\chi(k/K)\), where
\(\chi\) is even and smooth, supported in \([-1,1]\) and
equal to one on \([-1/2,1/2]\). Use the product filter in all
angles. Its kernel is the periodization of a rescaled Schwartz
function on Euclidean space, by Fourier inversion, and it
reproduces every trigonometric polynomial whose coordinate
degrees are at most \(K/2\).
For \(G(z)=\psi(z)-z\log r_*\), applying this filter to
\(G(c_S(x))\) as a function of angles produces a polynomial
\(L_{S,x}\) in the normalized means, of degree \(O(K)\),
with the following bound at the true angles on valid intervals
for any fixed \(b'>0\):
\begin{equation}\label{eq:filter}
 |L_{S,x}-G(c_S(x))|\le
        C_{b'} D_S(1+|c_S(x)|/D_S)^{-b'}.
\end{equation}
To prove the bound, write \(c=c_S(x)\) at the true means
and \(v\) for the increment due to the angle shift.
With \(z/K\) denoting the Euclidean shift before
periodization, \eqref{eq:increment} bounds \(|v|\) by
\(CD_S(|z|+|z|^2)\).
If \(|c|\le D_S\), subtract the exactly reproduced linear term
\(z'\log(D_S/r_*)\) in the function argument \(z'\), and
rescale the remaining \(z'\log|z'/D_S|\) by \(D_S\).
The Schwartz bound gives error \(O(D_S)\). If \(|c|>D_S\),
subtract the Taylor polynomial in \(v\) at \(c\) up to some
fixed order \(a-1\), with \(a\ge2\). Its powers of the increment
have finite angle degree at most \(a-1\) and are reproduced
for large \(K\). On \(|v|\le |c|/2\) the error is bounded by
\(C_a|v|^a|c|^{1-a}\). On the complement,
\(|z|+|z|^2\ge c'|c|/D_S\). After cancellation of the linear
part containing \(\log(|c|/r_*)\), the remaining difference,
including the Taylor terms, is bounded by \(|c|\) times a fixed
polynomial in \(1+|v|/|c|\). This follows by scaling
\(\psi(c+v)\) by \(|c|\) and bounding
\(|u\log|u||\le C(1+|u|^2)\). Choosing \(a\) and the
Schwartz decay order large proves \eqref{eq:filter}.
The approximation is even in each angle, which justifies
the polynomial representation.

On all rectangle arguments, \(|c_S(x)|\le C a_*\) in
either model. The Fourier coefficients of \(G(c_S(x))\)
are then at most \(C a_*(1+|\log(r_*/a_*)|)\) in
absolute value. As for the selection probabilities, conversion
from angle modes gives coefficient norm at most
\(C^K a_*(1+|\log(r_*/a_*)|)\), uniformly in \(x\).

\subsection{The recursive partition and its bias}

At stage \(j<J\), use the positive selection at \(S_j\) to
produce the rounded mass \(\widehat p_j\) and the rounded
projection \(a_j\), for which \(a_j/\ell_j\) uses the linear
form for \(c_j=c_{S_j}\). Set
\[
 B_j=L_2 r_*(1+\widehat p_j/r_*)^\alpha,\qquad
 T_j=\{x\in X_j:|a_j(x)|>\ell_j(x)B_j\},\qquad
 S_{j+1}=X_j\setminus T_j .
\]
Use \(X_j=S_j\) for \(j\ge1\), take \(T_J=S_J\),
and let null regions have zero contribution.
All possible cuts from any rounded vector in the fixed
bounded mesh range give a finite dictionary before observing
the interval-pilot sample. Its cardinality is polynomial in \(n\), since
the dimension and \(J\) are fixed. Define the regional kernels
there in advance, independently of the pilot and correction
observations. The dictionary can depend on \(n,N\).
On \(T_j\), the identity
\(\ell_j G(q/\ell_j)+q\log(\ell_j r_*/a_*)=\psi(q)-q\log a_*\)
decomposes the target using the envelope \(\ell_j\) of the
current kernel. We approximate the \(G\) term at \(c_j=q_{S_j}/\ell_j\)
and account for its difference from \(q/\ell_j\) below.
The resulting polynomial contribution on \(T_j\) is
\begin{equation}\label{eq:piece}
 \int_{T_j}\ell_j(x)L_{S_j,x}\dd x+
 \int_{T_j}q(x)\log(\ell_j(x)r_*/a_*)\dd x .
\end{equation}
The last term is linear in the ordinary global score means.
For the weighted envelope, \(\ell_S\ge1\) and
\[
 \int_S\log\ell_S
   = \int_S\log B_S+v_S\log(1/v_S)
   \le C v_S+v_S\log(1/v_S)\le C.
\]
This also gives a uniform bound for its coefficient norm after
including \(\log(r_*/a_*)\). The coefficient functions of the observed endpoints and the
angles are Borel. In particular the bin indicator uses the
specified rounding map, while \(G\) and the filter are continuous.
The subsequent operations are Fourier integrals over the angle
torus and spatial integrals of Borel integrands on predetermined
regions. These define Borel coefficients for the polynomial in
the score means once the real pilot has been observed.

To check the bias, condition the artificial selections on the
preceding rounded vectors, and put \(u_j=p_{S_j}/r_*\).
At each stage the probability that the scaled projection or mass
moves by more than \(C r_*(1+u_j)^\alpha\) is at most
\(C_M(1+u_j)^{-M}\), for any prescribed fixed \(M\).
In fact, \eqref{eq:increment} shows that for large \(u_j\)
such a move requires \(K|z|\ge c(1+u_j)^{(\alpha-1/2)/2}\)
in some angle, and the positive kernel has tails of arbitrarily
large fixed order by choice of \(b\). Rounding contributes
\(O(r_*)\). Bounded \(u_j\) are included by increasing \(C_M\).
Call the complementary event a good selection.
On it \(1+\widehat p_j/r_*\asymp1+u_j\) by \(\alpha<1\).
Choosing \(L_2\) sufficiently large, a good selection gives
\[
 |c_j|\ge c r_*(1+u_j)^\alpha\ \text{on }T_j,\qquad
 |c_j|\le C r_*(1+u_j)^\alpha\ \text{on } S_{j+1}.
\]
Consequently \eqref{eq:filter}, with \(b'\) sufficiently large,
and \(\int_{X_j}\ell_j\le C\) bound the mean filter error at
each nonterminal stage by \(C r_*\). The bad selection term
uses \(C D_{S_j}(1+u_j)^{-M}\le C r_*\).

Also on a good selection,
\[
 p_{S_{j+1}}\le C r_*(1+u_j)^\alpha+Ce.
\]
This follows by integrating \(|q_{S_j}|\) over \(S_{j+1}\),
using the bound on \(|c_j|\) and \(\int_{X_j}\ell_j\le C\),
then using \eqref{eq:regional} and
\(\|w-q\|_{1,S_j}\le\sqrt{v_{S_j}}e\).
On other selections \(p_{S_{j+1}}\le p_{S_j}\).
For any \(0<b_0\le1\) these bounds give the conditional estimate
\[
 \E[(1+u_{j+1})^{b_0}\mid\text{preceding vectors}]
 \le C\bigl((1+u_j)^{\alpha b_0}+(1+e/r_*)^{b_0}\bigr).
\]
The terminal filter error is bounded by \(C\E D_{S_J}\),
where \(D_{S_J}=r_*(1+\sqrt{u_J})\).
Iterate the conditional bound starting with \(b_0=1/2\),
reducing the exponent by \(\alpha\) at each prior stage.
After \(J\) steps the exponent on \(1+u_0\) is \(\epsilon\).
All residual powers have exponent at most one, so each
contributes \(O(r_*+e)\) after multiplication by \(r_*\).
The terminal error is at most \(C[r_*(1+u_0^\epsilon)+e]\). Terminated paths
can be assigned zero mass and zero contributions.
If \(\inf_W|q|\ge p_0/2\), for sufficiently large \(u_0\)
the first good selection empties \(S_1\) since
\(r_*(1+u_0)^\alpha=o(p_0)\). On the bad first selection
use \(D_{S_J}\le D_{S_0}\), together with its small
probability. Bounded \(u_0\) give the same \(O(r_*)\) bound
directly. Thus the terminal power loss is confined to
\(\inf_W|q|<p_0/2\).

It remains to replace \(q_{S_j}\) by \(q\) in the nonlinear
term of \eqref{eq:piece}. On an interval
\([-C p_0,C p_0]\), replace \(G\) by zero on
\([-r_*,r_*]\) and leave it equal to \(G\) outside.
The resulting function is continuous with Lipschitz constant
\(C(1+\log(1+u_0))\) and differs from \(G\) by at most \(C r_*\).
Both \(q_{S_j}/\ell_j\) and \(q/\ell_j\) lie in that interval.
After multiplication by \(\ell_j\) and integration on \(T_j\),
\eqref{eq:regional} bounds the error of the replacement by
\[
 C r_*+ C(1+\log(1+u_0))\sqrt{|S_j|}\,e .
\]
At the ordinary root \(q_0=q\). On a good first selection,
\(|q|\le C r_*(1+u_0)^\alpha\) on \(S_1\); for large \(u_0\)
apply Lemma~\ref{lem:sublevel} to \(q/p_0\), of mean one, to
bound \(|S_1|\le C(1+u_0)^{-(1-\alpha)\beta}\).
On a bad first selection the probability bound suffices, since
\(|S_j|\le1\). The expected replacement error is therefore
\(O(r_*+e)\). Bounded \(u_0\), including \(p_0=0\),
need only the volume bound. The rescaling identity used to define \eqref{eq:piece} now
applies on every piece of the partition. Summing the mean
errors proves \eqref{eq:angular-bias}.

\subsection{Estimating the path polynomial}

Average the contributions \eqref{eq:piece} exactly with their
polynomial path probabilities, multiplying by the probability
of each selection up to and including the one defining \(T_j\).
No selections are sampled in the implemented estimate.
At each node the outgoing coefficient norms sum to at most
\(C^K\); submultiplicativity and the fixed depth give total
degree at most \(C_J K+1\) and total coefficient norm at most
\begin{equation}\label{eq:path-C}
 C_J^K a_*(1+|\log(r_*/a_*)|).
\end{equation}
This bounds the sum of coefficients over the entire path
expansion, including products of scores on different regions.

To estimate a monomial, use the ordered distinct products
of the actual normalized observations \(Y=(g-l)/B\).
In the fixed-count correction fraction of size \(m_c\), divide
the degree-\(k\) ordered sum by \((m_c)_k=m_c(m_c-1)
\cdots(m_c-k+1)\). It is unbiased, conditional on the real
pilot. On valid intervals \(|\E Y|\le1\) and
\(\E Y^2\le1+C/t\), since \(\Var(g)\le u\) and
\(B\asymp\sqrt{tu}+t\).
In the Poisson model normalize \(Y=g/A\).
If \(p_0\le A<1\), a single parent observation may be represented
by the law on \(\Omega\) with density \(w\) and extra mass
\(1-p_0\) at one point outside \(\Omega\), where all scores
are zero. Consequently \(\E Y\le1\) and \(\E Y^2\le1/A\),
using \(g^2\le g\) pointwise. The Poisson ordered sum is divided
by \(m^k\), where \(m\) is the deterministic parent mean.
It can use just the restricted cloud since tuples with an
outside score vanish. The expectation is the product of the
parent score means, by the factorial moments of \(\Pois(m)\).
Equivalently, the local count law is \(\Pois(m p_0)\), and for
\(p_0>0\) its marks have conditional density \(w/p_0\), yielding
the same expectation. The degree-zero term is one in either
model, including with an empty cloud.

Here is a second moment bound that also applies to mixed scores.
Let \(B_2\) bound their individual second moments. Two ordered
\(k\)-tuples sharing \(i\) observations have pair contribution
at most \(B_2^i\) by Cauchy--Schwarz on each common observation
and the bound on unshared means. In the fixed-count group,
the pairing count after division by the squared denominator
is at most \(\binom{k}{i}(2k/m_c)^i\) for \(k\le m_c/2\).
For the second tuple, choose the positions of the common points,
assign distinct points from the first tuple, and bound the
remaining choices by \((m_c-i)_{k-i}\). Division by
\((m_c)_k\) yields the bound. The Poisson count instead gives
the factor \(m^{2k-i}/m^{2k}\) for the distinct points,
so the bound holds with \(2k/m\) in place of \(2k/m_c\).
Thus the second moment is at most \((1+2k B_2/m_c)^k\)
or its Poisson analogue. In our applications
\(mA\asymp\log n\) and \(m_c t\asymp\log n\).
Choose \(c_K\) so \(C_J K+1\le m_c/2\) at the
smallest allowed \(N\). The bounds on the monomials and
Minkowski with \eqref{eq:path-C} give second moments
\(C_J^K A^2\log^2 n\) and \(C_J^K\log^2 n\), respectively,
with \(C_J\) increased if necessary.
For any prescribed \(\eta>0\), further decrease \(c_K\)
to bound the exponential factor by \(C n^\eta\).
Our estimator is finite on arbitrary allowed actual pilot
histories, because the interval widths are positive and the
computed paths use only the finite set of Borel scores.

\section{The conditional estimator at large counts}\label{sec:high}

The second moment just proved suffices when the mass of a cell
is small. At larger counts we will obtain a conditional variance
bound \(C n^\eta\log^3(n)/N\) on the good pilot event.
Fix \(\gamma>0\) and suppose
\(n^\gamma\le N\le n\) with \(w=w_I\) a probability density.
We keep \(t=\log n/N\), \(q=q_I\), \(e=\|w-q\|_2\) and
\(r_*=t/K^2\). Split the observations into three independent
fractions of sizes comparable to \(N\), for a coarse pilot,
the interval pilot, and the correction group.

The empirical ordinary projection from the coarse pilot,
with coefficients rounded at scale \(n^{-3}\), gives a
polynomial \(a\) of uniformly bounded coefficients. It satisfies
\begin{equation}\label{eq:coarse}
 \|a-q\|_\infty\le \delta=L_3\sqrt{\log n/N}
\end{equation}
with probability at least \(1-n^{-M}\) for any prescribed
\(M\) with \(L_3\) sufficiently large, by concentration
of the fixed bounded basis scores. For a fixed \(D_0>2\), set
\[
 S=\{x\in W: |a(x)|\le D_0\delta\},\qquad T=W\setminus S .
\]
Conditional on the coarse pilot, take intervals for the
ordinary scores and the finite dictionary of kernels obtained
by starting the cuts of Section~\ref{sec:angular} at \(S\).
This includes the indicator of \(S\).
The real interval pilot is independent of \(a\) and the
correction sample. On \eqref{eq:coarse},
\(\sup_S|q|\le(D_0+1)\delta\) and \(|S|\le C\delta^\beta\),
by Lemma~\ref{lem:sublevel}. If \(\inf_W|q|\ge1/2\),
then \(S\) is empty for \(n\) large enough, uniformly in
\(n^\gamma\le N\le n\).

On \(S\) include the linear term \(\int_S q\log\delta\).
Include the nonlinear correction from the angular construction
only if \(|S|>0\) and the interval pilot's mean
\(\widehat p_S\) of \(\one_S\) satisfies
\(\widehat p_S\le L_4\delta\). On valid pilots, retention implies
\(p_S\le C\delta\), and rejection of a positive-volume \(S\)
implies \(p_S\ge (L_4/2)\delta\) for sufficiently large \(n\).
Both implications follow from
\(|\widehat p_S-p_S|\le C(\sqrt{tp_S}+t)\) with \(t=o(\delta)\).
Choose \(L_4/2>D_0+2\).
In the rejected case,
\[
 \sqrt{|S|}\,e\ \ge\ \int_S(w-q)
   \ \ge\ p_S-(D_0+1)\delta |S|\ \ge\ c\delta .
\]
Omitting the centered nonlinear term then incurs bias at
most \(C\delta|S|\le Ce\), since \(|q|\le(D_0+1)\delta\)
on \(S\). A null set contributes zero.

For a retained \(S\), start with both support and evaluation
region equal to \(S\), with its weighted kernel and initial
mass \(p_S\). In \eqref{eq:piece} the ordinary projection
\(q\) in the linear term still uses all of \(\Omega\).
Use \(a_*=\delta\), \(r_*=t/K^2\) and the same Bernstein
intervals throughout the \(J\) selections. The stage bias and mass
recurrence from Section~\ref{sec:angular} give
\(C r_*(1+(p_S/r_*)^\epsilon)+Ce\) for the filter.
The replacement of regional projections by \(q\) is needed
from the first weighted stage. At each stage the estimate is
\(C r_*+C\log n\sqrt{|S|}\,e\), by \eqref{eq:regional},
\(\|q\|_\infty\le C\), and \(r_*\gtrsim (n\log n)^{-1}\).
It is \(O(r_*+e)\) because \(|S|\le C\delta^\beta\),
whose positive power dominates the logarithm for
\(N\ge n^\gamma\). Thus the centered angular part, with the
gate included, has absolute conditional bias at most
\begin{equation}\label{eq:high-bias}
 C\bigl[e+r_*(1+r_*^{-\epsilon}
                        \one_{\{\inf_W|q|<1/2\}})\bigr],
\end{equation}
using \(p_S\le1\) and the empty-set assertion.

On valid retained histories, all descendant restricted-score
intervals have endpoints \(O(\delta)\), since their means
are bounded by \(p_S\le C\delta\) and their widths are
\(O(\sqrt{t p_S}+t)\). Thus the signed filter acts on
\(G(c)\) for \(|c|\le C\delta\) on the rectangles, giving
coefficient norm at most \(C^K\delta\log n\).
The global score intervals have widths \(O(\delta)\).
The ordinary projection evaluated at their lower endpoints
differs from \(q\) by \(O(\delta)\) uniformly, by the bounded
basis and fixed dimension. Its value on \(S\) is \(O(\delta)\).
Thus the constant term and each coefficient of that
projection as an affine expression in normalized scores is
\(O(\delta)\) on \(S\). Integrating its product with
\(\log(\ell_j r_*/\delta)\) on any descendant preserves a
coefficient norm bound \(C\delta\log n\), by the log-envelope
integral in Section~\ref{sec:angular}. Sum over the finite
paths via the outgoing probability norm bound. The factorial
moment argument now bounds the conditional second moment of
the centered angular estimate by \(C_J^K\delta^2\log^2 n\).
Its absolute conditional mean is at most \(C\delta\log n\):
at the true rectangle coordinates, the absolute filtered
value is bounded by \(C\delta\log n\) by the uniform
\(L^1\) norm of the filter, and the path weights form a
probability distribution with fixed depth.

On \(T\), put \(z=(q-a)/a\). On \eqref{eq:coarse} we have
\(|z|\le 1/D_0\). The pointwise identity
\[
 \psi(q)=q\log|a|+a(1+z)\log(1+z)
\]
allows the use of the Taylor polynomial of
\((1+z)\log(1+z)\) at zero of degree
\(M_T=\lceil C_T\log n\rceil\). Its nonconstant coefficients
have magnitude at most one. The integrated error is bounded by
\(C D_0^{-M_T}\int_T |a|\), which is \(o(r_*)\) for fixed
large \(C_T,D_0\). Estimate each power of \(z(x)\) using
ordered products of the score
\[
 Y(x,X)=\frac{K_\Omega(x,X)-a(x)}{a(x)}
\]
on the correction sample, with its actual falling factorial
denominator. Since \(\E_w K_\Omega(x,X)=q(x)\), the mean
of this signed score is \(z(x)\). On \(T\), the denominator
\(a(x)\) is separated
from zero by the observed threshold on every pilot history.
On \eqref{eq:coarse}, \(|\E Y|\le1/D_0\) and
\(\E Y^2\le C(1+N/(D_0^2\log n))\).
The overlap count from Section~\ref{sec:angular} now retains the
unshared means. For an order-\(k\) statistic, also with one
observation omitted, it bounds the second moment by
\[
 \bigl(D_0^{-2}+Ck(1+N/(D_0^2\log n))/N\bigr)^k .
\]
For \(k\le M_T\) this is at most \(\rho^{2k}\) with
a fixed \(\rho<1\), by choosing \(D_0\) sufficiently large
after \(C_T\), and then \(n\) large. The correction size
is of order \(N\) and \(M_T=o(N)\).

To bound the variance of the integrated Taylor statistic,
replace one correction observation \(X_i\) by an independent
copy \(X_i'\). At order \(k\) the difference of the ordered products statistic equals
\((k/m_c)(Y(x,X_i)-Y(x,X_i'))\) times the order-\((k-1)\)
statistic omitting \(i\). Here \(k/m_c=k(m_c-1)_{k-1}/(m_c)_k\).
The factors are independent conditional on the pilots, and
\[
 a(x)\bigl(Y(x,X_i)-Y(x,X_i')\bigr)
  = K_\Omega(x,X_i)-K_\Omega(x,X_i') .
\]
Thus after multiplication by \(a(x)\) the first factor has
\(L^2\) norm at most \(C k/m_c\)
by the fixed kernel bound. The other factor has \(L^2\) norm
at most \(\rho^{k-1}\). Minkowski for the integral and
Taylor sum gives \(O(1/m_c)\) for the \(L^2\) effect of
each replacement. The Efron--Stein variance bound (the sum
of these squared replacement effects divided by two) gives
variance \(O(1/N)\).

Estimate the linear contribution \(\int_T q\log|a|+
\int_S q\log\delta\) by averaging the corresponding
kernel integral on the correction sample. The score is bounded
by \(C\log n\) from the threshold on \(T\) and the uniform
bound on \(a\). This term has conditional variance
\(O(\log^2 n/N)\). The three parts may share the correction
observations, so add their conditional standard deviations.
The resulting statistic estimates \(\int_W\psi(q)\) with
bias bounded by \eqref{eq:high-bias} and variance
\begin{equation}\label{eq:high-var}
 C n^\eta\log^3 n/N
\end{equation}
on the stated good actual pilots, after decreasing \(c_K\)
as necessary. Its conditional mean is bounded by \(C\log n\):
use the retained angular mean bound above, the uniform Taylor
mean bound from the geometric series at \(|z|\le 1/D_0\),
and the bounded linear score. The estimator is defined by
observed conditions, and the probability of failure of the
analytic pilot bounds is \(O(n^{-M})\) conditional on \(N\).

\section{The Poisson experiment and fixed-sample risk}\label{sec:risk}

Choose a dyadic \(h\asymp(n\log n)^{-1/(s+d)}\), and set
\(\tau=h^s\), \(v=(3h)^d\), so that \(v\tau\asymp(n\log n)^{-1}\).
Write \(a_s=s/(s+d)\) and fix \(\eta,\gamma>0\) with
\(\gamma+\eta<1-2a_s\). Choose \(J\) so that
\(\epsilon=(3/4)^J/2<\min(\omega,1/(s+d))\).

Start with a \(\Pois(n/2)\) sample and give each point an
independent uniform label among six groups and an independent
continuous order mark. The six groups are used for selection,
low-mass correction, shape estimation, estimation of the
enlarged-cell mass, estimation of the evaluation-cell mass,
and estimation of the log density scale. Each group has
mean \(m=n/12\).
A local group count is computed on \(J_I^{\rm phys}\);
apply the inverse of \(T_I\) only after this restriction.
The local count is Poisson of mean \(mP_I\) and, conditional
on that count, the positions in unit coordinates are iid
of density \(w_I\). Ordering by independent marks applies a
permutation independent of the positions. These laws follow
by splitting the original count according to its point
categories. For disjoint physical supports, the generating
function of the counts for labels and support indicators is
\(\prod_j\exp(\lambda_j(z_j-1))\) with the appropriate means
\(\lambda_j\), by averaging the conditional multinomial
generating function. Given the counts, the marked positions
are iid within each category and independent across categories.
Thus the entire marked restrictions on disjoint physical
supports are independent. All cell estimators below use their
own marked restriction only.

Use the selector group to choose the low rule when its
count in \(J_I^{\rm phys}\) is less than \(D_{\rm sel}\log n\).
The low rule is valid at \(P_I\le A=C_A\log n/n\), and the
other rule at \(P_I\ge C_H\log n/n\).
First choose \(C_H\) sufficiently large for the nonlow rule.
Then take \(D_{\rm sel}\) so that for
\(P_I<C_H\log n/n\) the selector mean is at most half the
threshold, and \(C_A\) so that for \(P_I>A\) that mean
is at least twice the threshold. Poisson tails give failure
probability \(O(n^{-M})\) for any required fixed \(M\) by
increasing these constants. The validity ranges overlap.
In the nonlow rule, the shape group count \(N\) obeys
\(N\asymp nP_I\) and \(C_*\log n\le N\le n\) on an
event of probability at least \(1-O(n^{-M})\) in that
range, by the same concentration and \(P_I\le1\).
Outside the observed window use zero for the shape statistic.

\subsection{Low mass}
The local Poisson intensity in unit coordinates is
\(m w(x)\dd x\), where \(w=v f\circ T_I=P_I w_I\)
for a cell of positive mass \(P_I\). Its projection and residual
are \(q=P_I q_I\) and \(e=P_I E_I\).
For a point from the full correction label, the position
on the torus has density \(f\). Map it to \(T_I^{-1}\) of
its location if in \(J_I^{\rm phys}\), and to the outside
point of Section~\ref{sec:angular} otherwise. This gives
precisely the parent observation probability and restriction
used there. Put \(r_*=A/K^2\). By
\(\psi(P_I q_I)=P_I\psi(q_I)+q\log P_I\) for \(P_I>0\),
\begin{equation}\label{eq:low-decomp}
 F_I=\int_W [\psi(q)-q\log A]
       +\left(\mu_I-\int_W q\right)\log(P_I/A)
       +\mu_I\log(A/v).
\end{equation}
At \(P_I=0\) all terms are interpreted as zero.
The first term is estimated by Section~\ref{sec:angular},
with residual \(P_I E_I\) and indicator
\(\one_{\{\inf_W |q|<P_I/2\}}=\nu_I\), also valid at \(P_I=0\)
under the zero-mass convention.
The final term is estimated linearly with the score \(\one_W\).

For the middle term, a degree-\(K\) polynomial \(P_K\)
approximates \(x\log x\) uniformly on \([0,1]\) to error
\(C/K^2\), with coefficient norm at most \(C^K\).
For a construction, put \(x=(1+\cos\theta)/2\).
The cosine coefficients of \(\log x\) at modes \(j\ge1\) are
\(2(-1)^{j-1}/j\), by the radial
logarithm expansion of \(\log|1+\rho e^{i\theta}|^2-\log 4\)
as \(\rho\) increases to one. The series converges in \(L^2\)
by square summability, and its pointwise limit off
\(\theta=\pi\) identifies the \(L^2\) limit. Multiplication by
\(x\) combines two adjacent coefficients and the central one,
giving \(O(j^{-3})\) for \(j\ge2\). The absolutely summable
series for \(x\log x\) therefore has tail \(O(K^{-2})\)
uniformly, and its truncation has the coefficient bound by
the Chebyshev recurrence. Now \(B_{K-1}(x)=(P_K(x)-P_K(0))/x\)
is a polynomial with
\(x|B_{K-1}(x)-\log x|\le C/K^2\) for \(x>0\).
The multiplier \(\mu_I-\int_W q\) is the mean of
\[
 b(y)=\one_W(y)-\int_W K_\Omega(x,y)\dd x
\]
on the parent law, extended by zero outside \(\Omega\).
Its absolute value is \(O(P_I)\).
For clarity, the score \(b\) is represented in the statistic
as \(2D_b g_b-D_b\one_\Omega\), where
\(D_b\ge\max(1,\|b\|_\infty)\) and
\(g_b=\one_\Omega(1+b/D_b)/2\) belongs pointwise to \([0,1]\).
Use the factorial polynomial estimator for the product of
this mean and \(B_{K-1}(P_I/A)\).
Together with the linear term, its coefficient norm in the
scaled score means is \(O(C^K A\log n)\).
The bias from the replacement of the logarithm is at most
\(C A/K^2\), including zero mass by continuity of the
mass-multiplied target and the polynomial formula.
All terms have the parent denominator \(m^k\).
The factorial bound and Minkowski give a combined second
moment at most \(C n^\eta A^2\log^2 n\); decrease \(c_K\) if
needed for this combined bound. The bias is at most
\begin{equation}\label{eq:low-bias}
 C\bigl[P_I E_I+r_*(1+(P_I/r_*)^\epsilon
                          \nu_I)\bigr].
\end{equation}

\subsection{Nonlow mass}
Conditional on \(N\) in the observed range, the ordered
shape sample supplies the independent fractions in
Section~\ref{sec:angular} for \(N<n^\gamma\) and
Section~\ref{sec:high} for \(N\ge n^\gamma\).
The shape estimate of \(\int_W\psi(q_I)\) is multiplied
by the independent shape-mass count divided by \(m\), an
unbiased estimate \(\widehat P_I\) with variance \(P_I/m\).
On the good counts and good actual pilots, the resulting
conditional bias for \(P_I\int_W\psi(q_I)\) is bounded by
\begin{equation}\label{eq:mass-shape-bias}
 C\left[P_I E_I+\frac{P_I}{N\log n}
       \bigl(1+(N\log n)^\epsilon\nu_I\bigr)\right].
\end{equation}
Here \(r_*=\log n/(NK^2)\asymp (N\log n)^{-1}\).
Its conditional variance at these histories is
\(C n^\eta\log^C n\, P_I^2\) for \(N<n^\gamma\), and
\(C n^\eta\log^C n\,P_I/n\) for \(N\ge n^\gamma\).
Indeed \(\E\widehat P_I^2=O(P_I^2)\) on the valid
range. The product variance for an independent mass factor
and shape factor is this second moment times the shape
variance, plus the mass variance times the squared shape
mean. In the high range use \(N\asymp nP_I\) and the
\(C\log n\) mean bound from Section~\ref{sec:high}.
In the intermediate range the shape second moment alone
and \(nP_I\gtrsim\log n\) suffice.

Use two more independent groups for \(\mu_I\log b_I\).
Estimate \(\mu_I\) by the evaluation-cell count divided by
\(m\), and use the enlarged-cell count \(C_I\) in the
other group to form
\[
 L_{C_I}=\sum_{j=1}^{C_I}j^{-1}-\gamma_{\rm E}-\log(mv),
\]
where \(\gamma_{\rm E}\) is Euler's constant and the empty
sum is zero. For \(\lambda=mP_I\), the mean of the harmonic
sum equals \(\int_0^\lambda(1-e^{-u})\dd u/u\), by using
\(\sum_{j=1}^k j^{-1}=\int_0^1 (1-t^k)\dd t/(1-t)\)
and the Poisson generating function. For \(C_I\ge\lambda/2\) and \(\lambda\ge2\),
the harmonic sum differs from the sum at \(\lfloor\lambda\rfloor\)
by at most \(C|C_I-\lfloor\lambda\rfloor|/\lambda\).
The complementary event has exponentially small probability
and squared difference at most \(C(1+\log\lambda)^2\).
The expected squared difference is therefore at most \(C/\lambda\).
This bounds the variance and shows that the mean harmonic sum
minus \(\log\lambda\) tends to \(\gamma_{\rm E}\) as
\(\lambda\to\infty\). Subtracting \(\log\lambda\) in its
integral representation and integrating the derivative
\(-e^{-\lambda}/\lambda\) to infinity gives
\(\gamma_{\rm E}+\int_\lambda^\infty e^{-u}\dd u/u\).
As \(\lambda/(mv)=b_I\), the bias of \(L_{C_I}\) for \(\log b_I\)
is \(\int_\lambda^\infty e^{-u}\dd u/u\).
The bias can be made \(O(n^{-M})\) on the nonlow range
by \(C_H\) large. Independence, \(\mu_I\le P_I\) and
\(|\log b_I|\le C\log n\) there give variance at most
\(C P_I\log^2 n/n\) for the product with the estimate of
\(\mu_I\).

\subsection{Summation and return to fixed sample size}

Let \(Z_I\) be the output selected in cell \(I\), clipped
to \([-n^{B_0},n^{B_0}]\) for a fixed sufficiently large
\(B_0\). Let \(\mathcal H_I\) be the sigma-field of the
local selector and, on nonlow routes, the count \(N\) and
actual shape pilots used by the rule. The good event consists
of a valid route and, for a nonlow rule, count comparability
and validity of the real pilots. It is \(\mathcal H_I\)-measurable
for fixed \(f\), with failure probability \(O(n^{-M})\).
The estimator only checks observed conditions. In particular
it uses the actual interval-pilot regional count for the
gate in Section~\ref{sec:high}, and the auxiliary path
probabilities there and in Section~\ref{sec:angular} are
built into the estimator.

The good-history second moments of the unclipped outputs
are bounded by a fixed power of \(n\) uniformly over \(f,I\).
This follows from the bounds on their biases, variance and
means, with \(|F_I|\le C\log n\) even at small mass since
\(P_I|\log P_I|\) is bounded. One may choose the clipping
exponent before making the failure probability arbitrarily
small: the powers in the good-history bounds are controlled
by the fixed \(\eta,\gamma,s,d\), with arbitrary concentration
constants absorbed by enlarging the sample-size cutoff.
Clipping changes each good-history mean by \(o(v\tau)\)
using the second moment divided by the clipping threshold,
and contracts the conditional variance by the independent-copy
formula. On bad histories the squared conditional bias and
the conditional variance are bounded by
\(C(n^{B_0}+\log n)^2\). Choose \(M>2B_0+3\).
The sum of the bad-history conditional bias contributions
in \(L^2\) is at most
\(C h^{-d}(n^{B_0}+\log n)n^{-M/2}=o(\tau)\).
For the variance sum the bound is
\(C h^{-d}n^{2B_0-M}=o(\tau^2)\).

Combining \eqref{eq:low-bias} and
\eqref{eq:mass-shape-bias}, on all good histories
\begin{equation}\label{eq:per-cell}
 \left|\E[Z_I\mid\mathcal H_I]-F_I\right|
 \le C P_I E_I+ C v\tau\left(1+(b_I/\tau)^\epsilon
                                 \nu_I\right).
\end{equation}
On low routes \(A/K^2\asymp v\tau\) and
\(P_I/(A/K^2)\asymp b_I/\tau\) when \(P_I>0\).
On the good nonlow histories \(N\asymp nP_I\) gives
\(P_I/(N\log n)\asymp v\tau\) and \(N\log n\asymp b_I/\tau\).
The residuals sum to \(O(\tau)\) by \eqref{eq:projection}.
The terms for \(b_I\le\tau\) and the constant term
also sum to \(O(\tau)\). For the remaining term partition
the amplitudes into disjoint dyadic bands
\(B/2<b_I\le B\) with \(B=2^k\tau\), \(k\ge1\).
At \(\xi=B/\tau\), Lemma~\ref{lem:bands} gives
a contribution at most
\[
 C\tau\,\xi^\epsilon(\xi^{-\omega}+h)
       +C \xi^{\epsilon-1}\sum_{B/2<b_I\le B}\sqrt v\,e_I .
\]
The largest \(\xi\) is at most \(C h^{-(s+d)}\)
since \(P_I\le1\). The dyadic sum of \(\xi^{\epsilon-\omega}\)
is bounded since \(\epsilon<\omega\), and that of
\(h\xi^\epsilon\) is \(O(h^{1-(s+d)\epsilon})\). The residuals remain in disjoint
bands and \(\xi^{\epsilon-1}\le1\). Therefore the sum of the
conditional biases has \(L^2\) norm \(O(\tau)\) by Minkowski,
including the bad-history contribution.

For the centered part, color the intersection graph of the
physical supports by at most \(D_{\rm col}=7^d\) colors,
using the degree bound from Section~\ref{sec:geometry}.
Within each color the local marked restrictions are independent.
Both \(Z_I\) and its conditional mean given \(\mathcal H_I\)
are measurable with respect to this local restriction.
Consequently the centered terms have mean zero and are
independent inside each color. By Cauchy--Schwarz across colors,
\[
 \E\left|\sum_I (Z_I-\E[Z_I\mid\mathcal H_I])\right|^2
 \le D_{\rm col}\sum_I \E\Var(Z_I\mid\mathcal H_I).
\]
On good low histories the sum of the conditional variances is
bounded by \(C h^{-d} n^\eta A^2\log^2 n\). Good intermediate
histories occur only for \(P_I\le C n^{-1+\gamma}\),
by \(N<n^\gamma\) and count comparability there.
The good intermediate variances thus sum, in expectation,
to at most \(C n^{-1+\gamma+\eta}\log^C n\),
using \(\sum_I P_I\le5^d\).
The good high histories contribute
\(C n^{-1+\eta}\log^C n\) by the same mass sum.
Together these bounds give
\[
 C\log^C n\bigl(h^{-d}n^\eta A^2+n^{-1+\gamma+\eta}\bigr)
   +o(\tau^2)=O(\tau^2).
\]
For the first term the strict power slack is
\(\eta<1-a_s\), and for the others it is
\(\eta+\gamma<1-2a_s\). Both follow from our choices.

The Poisson output \(\sum_I Z_I\) is bounded for each \(n\)
and measurable. The centered variance and conditional bias
bounds and \eqref{eq:transfer} give its root-mean-square
error at most \(C\tau\) for \(\int f\log f\), uniformly
over the class. The mesh cutoff \(h\le1/8\) and all
stochastic cutoffs are achieved by increasing a single
\(n_0(d,s,p,L)\), since all estimates above are uniform.

To use exactly \(n\) observations, draw \(M\sim\Pois(n/2)\)
independently. If \(M\le n\), apply the Poisson procedure to
the first \(M\) observations with independent auxiliary labels
and marks, and reverse its sign. Otherwise output zero.
Couple to an infinite iid sequence to see that the squared
error on \(M\le n\) is at most the squared error in the
complete Poisson experiment. On \(M>n\) the squared error
is \(H(f)^2\), bounded uniformly by \(\|f\|_2\le L\) and
\(|u\log u|\le C(1+u^2)\). The overflow probability is at most
\(\exp[-n(\log 2-1/2)]\) by the exponential moment at
\(\log 2\). The auxiliary-randomness average at each sample
vector is a finite measurable estimator. Jensen's inequality
gives no larger squared error for this deterministic estimator.
This proves the upper bound for \(2<s<d\). The likelihood
mixture of Section~\ref{sec:large-s} gives the rest of the
upper bound, completing Theorem~\ref{thm:main}.


\section*{AI use}
Codex assisted with the proof development, natural-language
checking, and preparation of the manuscript.

\begin{thebibliography}{3}
\bibitem{HJWW}
Y.~Han, J.~Jiao, T.~Weissman and Y.~Wu,
\textit{Optimal rates of entropy estimation over Lipschitz balls},
Ann. Statist. \textbf{48} (2020), 3228--3250.
\href{https://doi.org/10.1214/19-AOS1927}{doi:10.1214/19-AOS1927}.
\bibitem{HJWWv4}
Y.~Han, J.~Jiao, T.~Weissman and Y.~Wu,
\textit{Optimal rates of entropy estimation over Lipschitz balls},
author manuscript, version 4.
\href{https://arxiv.org/abs/1711.02141v4}{arXiv:1711.02141v4}.
\bibitem{YangBarron}
Y.~Yang and A.~Barron,
\textit{Information-theoretic determination of minimax rates of
convergence}, Ann. Statist. \textbf{27} (1999), 1564--1599.
\href{https://doi.org/10.1214/aos/1017939142}{doi:10.1214/aos/1017939142}.
\end{thebibliography}
\end{document}
